\begin{tikzpicture}[font=\small, >=Latex]
\tikzstyle{source}=[
    rectangle, rounded corners=2pt, draw=DTBlue, fill=DTLightBlue,
    text width=3.8cm, minimum height=1.0cm, align=center, font=\scriptsize
]
\tikzstyle{module}=[
    rectangle, rounded corners=2pt, draw=DTGreen, fill=DTLightGreen,
    text width=3.0cm, minimum height=0.78cm, align=center, font=\scriptsize
]
\tikzstyle{target}=[
    rectangle, rounded corners=2pt, draw=DTOrange, fill=DTLightOrange,
    text width=4.2cm, minimum height=1.0cm, align=center, font=\small
]

\node[source] (ehealth) at (0,3.0) {eHealth state of the art\\architecture, interoperability, governance};
\node[source] (cardio) at (0,1.4) {Cardiovascular state of the art\\biosignals, hemodynamics, patient specificity};
\node[source] (ai) at (0,-0.2) {AI and hybrid modeling\\diagnosis support, prediction, simulation};
\node[source] (val) at (0,-1.8) {Validation requirements\\uncertainty, traceability, human supervision};

\node[target] (ct) at (5.6,0.6) {CardioTwin PFE\\patient-oriented cardiovascular Digital Twin prototype};

\foreach \n in {ehealth,cardio,ai,val}{
    \draw[->, thick, DTBlue] (\n.east) -- (ct.west);
}

\node[module] (m1) at (10.5,3.2) {Biomedical data acquisition};
\node[module] (m2) at (10.5,2.1) {ECG preprocessing};
\node[module] (m3) at (10.5,1.0) {AI diagnostic assistance};
\node[module] (m4) at (10.5,-0.1) {Clinical and physiological assessment};
\node[module] (m5) at (10.5,-1.2) {Patient-state fusion and update};
\node[module] (m6) at (10.5,-2.3) {What-if simulation};
\node[module] (m7) at (10.5,-3.4) {Visualization and evaluation};

\foreach \n in {m1,m2,m3,m4,m5,m6,m7}{
    \draw[->, thick, DTGreen] (ct.east) -- (\n.west);
}
\end{tikzpicture}
